\documentclass[10pt,a4paper]{amsart}
\usepackage[margin=19mm]{geometry}
\usepackage{amsmath,amssymb,mathtools}
\usepackage[scr=rsfs,cal=euler]{mathalfa}
\usepackage{xfrac}
\usepackage[unicode,hidelinks,bookmarks=false]{hyperref}
\newcommand{\ie}{i.e.\ }
\newcommand{\cf}{cf.\ }
\newcommand{\resp}{respectively\ }
\newcommand{\Subsection}{\subsection}
\providecommand{\nobreakdash}{\nobreak\hbox{-}}
\setlength{\emergencystretch}{2em}
\multlinegap0pt
\usepackage{bm}
\usepackage{bbm}
\usepackage{dsfont}
\usepackage{xspace}
\usepackage{cancel}
\usepackage{mathtools}
\usepackage{amsthm}
\makeatletter
\@ifundefined{thm}{\newtheorem{thm}{Thm}}{}
\@ifundefined{prop}{\newtheorem{prop}[thm]{Proposition}}{}
\makeatother
\DeclareMathOperator{\Endo}{End}
\DeclareMathOperator{\uhp}{\mathcal{H}}
\providecommand{\pseries}[2]{#1[\![ #2 ]\!]}
\title[Vertex operator algebra bundles on modular curves and their ]{Vertex operator algebra bundles on modular curves and their associated modular forms}
\author{Daniel Barake, Owen Chuchman, Cameron Franc, Geoffrey Mason, Brett Nasserden}
\date{Source: arXiv:2601.10686v1 (CC BY 4.0)}
\begin{document}
\maketitle

\noindent\emph{Source and licence.} Vertex operator algebra bundles on modular curves and their associated modular forms. Version arXiv:2601.10686v1 (CC BY 4.0), distributed under the Creative Commons Attribution 4.0 International licence, as recorded in the arXiv licence metadata for this exact version and shown on the abstract page https://arxiv.org/abs/2601.10686v1. Full rights notice: licenses/arxiv-2601.10686-CC-BY-4.0.txt.\par\smallskip

\noindent\textit{Editorial scope.} Counted unit: the Prop whose source label is \texttt{p:positivemodes} in the article (source0.tex lines 1103--1112, source label \texttt{p:positivemodes}), delivered together with the article's own complete original proof of it (source0.tex lines 1113--1194, one \texttt{proof} environment of 82 lines). No mathematics is written, repaired or re-derived: statement and proof are the article's own text line for line. Required context retained uncounted: none. Every definition, display and label of those lines is retained unchanged. Omitted from this package and not redistributed: 1--1102, 1195--2635. Nothing inside a retained excerpt is omitted or altered: every retained body is delivered exactly as the article's lines are written, so no omission receipt of any kind accompanies this package. Citations: the retained excerpts carry no citation command at all, so no bibliography entry of the article is transcribed and no reference is redistributed. Reference adaptation: none was needed and none was made; every reference inside the delivered unit resolves inside it, so no unresolved reference or citation is produced. Printed slips kept literal: no printed word, symbol, hypothesis or number of the counted statement or of its proof was altered. Layout and class: the article's own class is \texttt{amsart} with packages including amssymb, latexsym, amsmath, amsthm, amscd, setspace, xy, geometry, graphicx, xcolor, mathtools, charter, hyperref, makecell; the wrapper is the lane's pinned \texttt{amsart} editorial wrapper at 10pt on A4, which restates the theorem-like environments the retained text needs and adds the article's own macro definitions that the retained text needs (\texttt{\textbackslash Endo}, \texttt{\textbackslash pseries}, \texttt{\textbackslash uhp}), with extra wrapper packages (bm, bbm, dsfont, xspace, cancel, mathtools). The delivered numbering is the unit's own. Assets: no retained span contains a figure environment, a graphics-inclusion command, a PDF-page inclusion, a table or a TikZ picture; no image file is copied into this package, the package declares no asset dependency and no image file is redistributed with it. Licence: version arXiv:2601.10686v1 of this article is distributed under the Creative Commons Attribution 4.0 International licence, as recorded in the arXiv licence metadata for this exact version and shown on the abstract page https://arxiv.org/abs/2601.10686v1. A full rights notice is delivered at licenses/arxiv-2601.10686-CC-BY-4.0.txt. Characters: every retained excerpt is pure ASCII, so the delivered engine, XeLaTeX with its default Latin Modern text font, reports no missing character.
\begin{prop}
\label{p:positivemodes}
    If $f$ and $g$ are VOA-valued modular forms, of weights $k$ and $\ell$ respectively, then for $m\geq -1$, the $m$th mode $f(\tau_1)(m)g(\tau_2)$ satisfies the transformation law
 \begin{align}
        \label{eq:prodtran}
        &f(\gamma\tau_1)(m)g(\gamma\tau_2)\\
        \nonumber =&j(\gamma,\tau_1)^kj(\gamma,\tau_2)^{\ell-2m-2}K(\gamma,\tau_2)\left(\sum_{n=0}^{\infty}\binom{n+m+1}{m+1}(-X(\gamma, \tau_2))^2f(\tau_1)(n+m)g(\tau_2)\right).
    \end{align}
    where $\tau_1,\tau_2 \in \uhp$ are distinct points and $\gamma \in \Gamma$.
\end{prop}

\begin{proof}
    First let $A\in V$ be any state. If we view $\tau_2$ as being fixed and $\tau_1$ as varying, then we have a local coordinate $\tau_1-\tau_2$ at $\tau_1=\tau_2$. For $\gamma \in \Gamma$, we then have two coordinates at $\tau_1 = \gamma \tau_2$: one is $\tau_1 - \gamma \tau_2$, and the other is $\gamma^{-1}\tau_1 - \tau_2$.  These coordinates differ by a change of variable, and we shall prove our theorem by comparing Huang's transformation formula with respect to these two local coordinates. Recall that his formula says the following: if $z = \tau_1-\gamma \tau_2$ then 
    \[
     Y(A,z) = R(\rho)Y(R(\rho_z)^{-1}A,\rho(z))R(\rho)^{-1}
    \]
    where $\gamma^{-1}\tau_1-\tau_2 = \rho(z)$. Recall that we saw earlier that
    \[
    \gamma^{-1}\tau_1-\tau_2 = \sum_{n\geq 1} c^{n-1}j(\gamma,\tau_2)^{n+1}(\tau_1-\gamma\tau_2)^n.
    \]
   and hence
    \[
    R(\rho) = e^{-cj(\gamma,\tau_2)L(1)}j(\gamma,\tau_2)^{-2L(0)}.
    \]
    We also need to determine $R(\rho_z)$ where $z = \tau_1-\gamma\tau_2$. Since we can write 
    \[\rho(X) = \gamma^{-1}(X+\gamma \tau_2)-\tau_2,\] 
    we find that
    \begin{align*}
        \rho_z(X) &=\rho(X+z)-\rho(z)\\
        &= \gamma^{-1}(X+z+\gamma \tau_2)-\tau_2 - (\gamma^{-1}(z+\gamma\tau_2) - \tau_2)\\
        &=\gamma^{-1}(X+\tau_1) - \gamma^{-1}\tau_1
    \end{align*}
    Therefore, by a previous computation setting $\tau = \gamma^{-1}\tau_1$, we have
    \[
        \rho_z(X) = \sum_{n\geq 1} c^{n-1}j(\gamma, \gamma^{-1}\tau_1)^{n+1}X^n.
    \]
    It follows that
    \[
      R(\rho_z) = e^{-cj(\gamma,\gamma^{-1}\tau_1)}j(\gamma,\gamma^{-1}\tau_1)^{-2L(0)}.
    \]
    Thus, in terms of our notation $K(\gamma,\tau) = e^{-cj(\gamma,\tau)}j(\gamma,\tau)^{-2L(0)}$, we have
    \begin{align*}
    R(\rho) &= K(\gamma, \tau_2),\\
    R(\rho_z) &= K(\gamma,\gamma^{-1}\tau_1).
    \end{align*}
    
    Putting all this together, Huang's formula is the identity:
    \begin{equation}
        \label{eq:huang1}
        Y(A,\tau_1-\gamma \tau_2) = K(\gamma, \tau_2)Y(K(\gamma,\gamma^{-1}\tau_1)^{-1}A,\gamma^{-1}\tau_1-\tau_2) K(\gamma,\tau_2)^{-1}.
    \end{equation}
    This holds for $\tau_1\neq \gamma \tau_2$.

    Recall that in terms of our new notation, VOA valued modular forms satisfy the transformation law
    \[
      f(\gamma \tau) = j(\gamma,\tau)^kK(\gamma,\tau)f(\tau).
    \]
    In particular, taking $\tau = \gamma^{-1}\tau_1$, this says
    \[
      f(\tau_1) = j(\gamma,\gamma^{-1}\tau_1)^kK(\gamma,\gamma^{-1}\tau_1)f(\gamma^{-1}\tau_1).
    \]
    Thus, if we set $A = f(\tau_1)$ in Huang's formula \eqref{eq:huang1}, where $f$ is a VOA-valued modular form, then we deduce that
        \begin{equation}
        \label{eq:huang2}
        Y(f(\tau_1),\tau_1-\gamma \tau_2) = j(\gamma,\gamma^{-1}\tau_1)^kK(\gamma, \tau_2)Y(f(\gamma^{-1}\tau_1),\gamma^{-1}\tau_1-\tau_2) K(\gamma,\tau_2)^{-1}.
    \end{equation}
    Now, if $g(\tau_2)$ is another VOA-valued modular form, we can multiply on the right by $g(\gamma \tau_2)$ to deduce that
        \begin{equation}
        \label{eq:huang3}
        Y(f(\tau_1),\tau_1-\gamma \tau_2)g(\gamma \tau_2) = j(\gamma, \gamma^{-1}\tau_1)^kj(\gamma,\tau_2)^{\ell}K(\gamma, \tau_2)Y(f(\gamma^{-1}\tau_1),\gamma^{-1}\tau_1-\tau_2) g(\tau_2).
    \end{equation}

    We now complete the proof by comparing the principal parts of both sides of equation \eqref{eq:huang3}. The extra weight-increasing automorphy factors arise by expressing $\gamma^{-1}\tau_1-\tau_2$ on the right hand side of equation \eqref{eq:huang3} in terms of $\tau_1-\gamma\tau_2$.

    Indeed, to simplify notation, if we set $B = f(\gamma^{-1}\tau_1)$, $E = \Endo(V)$, and $T=\tau_1-\gamma\tau_2$, then we find that
    \begin{align*}
    &Y(B,\gamma^{-1}\tau_1-\tau_2)\\
    \equiv& \sum_{n\geq -1} B(n) (\gamma^{-1}\tau_1-\tau_2)^{-n-1}\pmod{T\pseries{E}{T}}\\
    \equiv& \sum_{n\geq -1} B(n) \left(\sum_{m\geq 1} c^{m-1}j(\gamma,\tau_2)^{m+1}T^m\right)^{-n-1}\pmod{T\pseries{E}{T}}\\
    \equiv& \sum_{n\geq -1} j(\gamma,\tau_2)^{-2n-2}B(n)T^{-n-1} \left(\sum_{m\geq 1} \left(cj(\gamma,\tau_2)T\right)^{m-1}\right)^{-n-1}\pmod{T\pseries{E}{T}}\\
    \equiv& \sum_{n\geq -1} j(\gamma,\tau_2)^{-2n-2}B(n)T^{-n-1} \left(1-cj(\gamma,\tau_2)T\right)^{n+1}\pmod{T\pseries{E}{T}}\\
    \equiv& \sum_{n\geq -1}\sum_{m=0}^{n+1}(-1)^{n+1-m}\binom{n+1}{m}c^{n+1-m} j(\gamma,\tau_2)^{-m-n-1}B(n)T^{-m} \pmod{T\pseries{E}{T}}\\
    \equiv& \sum_{m=0}^\infty(-c)^{-m}j(\gamma,\tau_2)^{-m}\left(\sum_{n=m-1}^\infty\binom{n+1}{m}(-c)^{n+1} j(\gamma,\tau_2)^{-n-1}B(n)\right)T^{-m} \pmod{T\pseries{E}{T}}\\
    \equiv& \sum_{m=-1}^\infty j(\gamma,\tau_2)^{-2m-2}\left(\sum_{n=0}^\infty\binom{n+m+1}{m+1}(-c)^{n} j(\gamma,\tau_2)^{-n}B(n+m)\right)T^{-m-1} \pmod{T\pseries{E}{T}}
    \end{align*}
    Hence, combining this with equation \eqref{eq:huang3} we obtain for $m\geq -1$:
    \begin{align}
        \label{eq:prodtransformation}
        &f(\tau_1)(m)g(\gamma\tau_2)\\
        \nonumber =&j(\gamma,\gamma^{-1}\tau_1)^kj(\gamma,\tau_2)^{\ell-2m-2}K(\gamma,\tau_2)\left(\sum_{n=0}^{\infty}\binom{n+m+1}{m+1}(-c)^nj(\gamma,\tau_2)^{-n}f(\gamma^{-1}\tau_1)(n+m)g(\tau_2)\right).
    \end{align}
    Finally, if we replace $\tau_1$ by $\gamma \tau_1$, this proves the proposition.
    \end{proof}

\end{document}
